\booksection{ملحق الفصل 11: المعارف التقنية للتصوير الاحترافي}

\chapterintro{لفهم أسرار التصوير والفيديو بالهاتف حقاً، تحتاج إلى معرفة بعض المعلومات التقنية الاحترافية. هذه المعارف لا تساعدك فقط على استخدام كاميرا الهاتف بشكل أفضل، بل تجعلك أكثر قدرة على الحكم عند اختيار الأجهزة.}

\booksubsection{أولاً: الدقة ومعدل الإطارات ومعدل البت: العناصر الثلاثة لجودة الفيديو}

\subsubsection*{الدقة (Resolution):}

تحدد الدقة وضوح الصورة، وتُعبَّر عنها بعدد البكسلات.

\begin{itemize}
  \item \textbf{720P (HD)}: 1280×720 بكسل، دقة تمهيدية.
  \item \textbf{1080P (Full HD)}: 1920×1080 بكسل، الدقة الأكثر شيوعاً حالياً.
  \item \textbf{2K (QHD)}: 2560×1440 بكسل، أوضح من 1080P.
  \item \textbf{4K (UHD)}: 3840×2160 بكسل، فائق الوضوح، غني بالتفاصيل.
  \item \textbf{8K}: 7680×4320 بكسل، معيار المستقبل، ودعم الأجهزة له محدود حالياً.
\end{itemize}

كلما زادت الدقة، زاد وضوح الصورة، لكن الملف يكبر أيضاً، وتزداد المتطلبات على التخزين وقدرة المعالجة.

\subsubsection*{معدل الإطارات (Frame Rate):}

يحدد معدل الإطارات سلاسة الصورة، ويُعبَّر عنه بعدد الإطارات في الثانية (fps).

\begin{itemize}
  \item \textbf{24fps}: معدل إطارات الأفلام السينمائية القياسي، يتميز بـ "إحساس سينمائي" فريد.
  \item \textbf{30fps}: شائع في التلفزيون والفيديو على الإنترنت، سلس وطبيعي.
  \item \textbf{60fps}: معدل إطارات مرتفع، مناسب للمشاهد الرياضية وتسجيل الألعاب، سلس جداً.
  \item \textbf{120fps/240fps}: معدل إطارات فائق، يمكن استخدامه لفيديوهات الحركة البطيئة.
\end{itemize}

كلما زاد معدل الإطارات، زادت سلاسة الصورة وتحسّن تأثير الحركة البطيئة، لكن الملف يكبر أيضاً.

\subsubsection*{معدل البت (Bitrate):}

يحدد معدل البت كثافة بيانات الفيديو، ويُعبَّر عنه بحجم البيانات في الثانية (Mbps).

\begin{itemize}
  \item كلما زاد معدل البت، تحسّنت جودة الصورة وزادت التفاصيل المحفوظة، لكن الملف يكبر أيضاً.
  \item معدل البت المنخفض جداً يؤدي إلى ظهور فسيفساء وكتل لونية وآثار ضغط في الصورة.
  \item التصوير الاحترافي للفيديو يتطلب عادة معدل بت مرتفعاً للحفاظ على مساحة أكبر للمعالجة اللاحقة.
\end{itemize}

\booksubsection{ثانياً: صيغ ترميز الفيديو: H.264 وH.265 وProRes}

\subsubsection*{H.264 (AVC):}

\begin{itemize}
  \item \textbf{الميزات}: صيغة الترميز الأكثر استخداماً، توافقها ممتاز جداً.
  \item \textbf{المزايا}: يمكن لجميع الأجهزة تقريباً تشغيلها، وكفاءة الضغط متوسطة.
  \item \textbf{العيوب}: مقارنة بالصيغ الجديدة، كفاءة الضغط أقل، والملفات أكبر.
  \item \textbf{سيناريوهات الاستخدام}: مشاركة الفيديو على الإنترنت، التصوير اليومي للمستخدم العادي.
\end{itemize}

\subsubsection*{H.265 (HEVC):}

\begin{itemize}
  \item \textbf{الميزات}: النسخة المطوّرة من H.264، تحسّنت كفاءة الضغط بنحو 50\%.
  \item \textbf{المزايا}: ملفات أصغر عند الجودة نفسها، ويدعم دقة ومعدل إطارات أعلى.
  \item \textbf{العيوب}: الترميز وفك الترميز يتطلبان معالجاً أقوى، وبعض الأجهزة القديمة غير متوافقة.
  \item \textbf{سيناريوهات الاستخدام}: تصوير فيديو 4K، الحالات التي تكون فيها مساحة التخزين محدودة.
\end{itemize}

\subsubsection*{Apple ProRes:}

\begin{itemize}
  \item \textbf{الميزات}: صيغة ترميز فيديو احترافية طورتها Apple، ضغط شبه فقداني.
  \item \textbf{المزايا}: جودة صورة عالية جداً، مساحة كبيرة للمعالجة اللاحقة، ألوان دقيقة.
  \item \textbf{العيوب}: الملفات ضخمة جداً، تتطلب مساحة تخزين كبيرة، والمعالجة تتطلب معدات احترافية.
  \item \textbf{سيناريوهات الاستخدام}: إنتاج الفيديو الاحترافي، التصوير بمستوى سينمائي.
\end{itemize}

\subsubsection*{ProRes RAW:}

\begin{itemize}
  \item \textbf{الميزات}: يجمع بين مزايا ProRes وصيغة RAW.
  \item \textbf{ما هو RAW}: بيانات المستشعر الخام غير المعالجة، تحتفظ بأقصى قدر من المعلومات.
  \item \textbf{المزايا}: مساحة كبيرة جداً للتعديل اللاحق، يمكن تغيير التعريض وتوازن اللون الأبيض وغيرها دون فقدان الجودة.
  \item \textbf{العيوب}: الملفات ضخمة جداً، تتطلب برامج معالجة لاحقة احترافية.
  \item \textbf{سيناريوهات الاستخدام}: إنتاج احترافي راقٍ. سلسلة iPhone 15 Pro وما بعدها تدعم تصوير فيديو Apple ProRes (حتى 4K 60fps)، وiPhone 17 Pro يدعم بشكل إضافي تسجيل ProRes RAW مباشرة داخل الجهاز.
\end{itemize}

\booksubsection{ثالثاً: HDR: تقنية المدى الديناميكي العالي}

تقنية HDR (High Dynamic Range) تحافظ في الوقت نفسه على تفاصيل أكثر الأجزاء سطوعاً وأكثرها ظلاماً في الصورة.

\begin{itemize}
  \item \textbf{HDR10}: معيار HDR الأساسي، توافقه جيد.
  \item \textbf{HDR10+}: معيار تدعمه Samsung، يدعم البيانات الوصفية الديناميكية.
  \item \textbf{Dolby Vision}: معيار HDR من الطراز الأول طورته Dolby، جودته الأفضل، لكنه يتطلب ترخيصاً.
  \item \textbf{HLG}: معيار HDR بمستوى البث، متوافق مع الشاشات العادية.
\end{itemize}

عند تصوير فيديو HDR، تحتاج إلى شاشة تدعم HDR لرؤية التأثير الكامل.

\booksubsection{رابعاً: فضاء الألوان وعمق اللون}

\subsubsection*{فضاء الألوان:}

\begin{itemize}
  \item \textbf{sRGB}: فضاء الألوان القياسي، توافقه الأفضل.
  \item \textbf{Display P3}: معيار نطاق لوني واسع تروّج له Apple، ألوانه أغنى.
  \item \textbf{Rec.2020}: فضاء الألوان القياسي لفيديو HDR، نطاقه اللوني الأوسع.
\end{itemize}

\subsubsection*{عمق اللون:}

\begin{itemize}
  \item \textbf{8-bit}: 16.7 مليون لون، معيار الفيديو العادي.
  \item \textbf{10-bit}: 1.07 مليار لون، معيار الفيديو الاحترافي، انتقالات لونية أكثر سلاسة.
  \item \textbf{12-bit}: 68 مليار لون، معيار احترافي راقٍ.
\end{itemize}

\begin{tcbinfo}{توصيات عملية}
\begin{itemize}
  \item التصوير اليومي: 1080P/4K + 30fps + ترميز H.265، توازن بين الجودة وحجم الملف.
  \item تصوير الحركة: 1080P/4K + 60fps، لضمان السلاسة.
  \item الإبداع الاحترافي: 4K + ProRes (إذا كان مدعوماً)، للحفاظ على أقصى مساحة للمعالجة اللاحقة.
  \item توصية التخزين: الفيديو بمعدل بت مرتفع يستهلك مساحة كبيرة، يُنصح بالنسخ الاحتياطي الدوري إلى الكمبيوتر أو التخزين السحابي.
\end{itemize}

\end{tcbinfo}

\begin{tcbpractice}{استكشف: إعدادات الفيديو في هاتفك}
\begin{enumerate}
  \item افتح تطبيق الكاميرا، وبدّل إلى وضع الفيديو.
  \item اعرض خيارات الدقة ومعدل الإطارات، وجرّب مجموعات مختلفة.
  \item إذا كان مدعوماً، تحقق من وجود خيارات متقدمة مثل HDR وProRes.
  \item صوّر عدة مقاطع بإعدادات مختلفة، وقارن الجودة وحجم الملف.
\end{enumerate}

\end{tcbpractice}

\booksubsection{خامساً: مقارنة صيغة RAW مع HEIF/JPEG}

\subsubsection*{ما هي صيغة RAW:}

صيغة RAW هي “البيانات الخام” التي يسجلها مستشعر الكاميرا — معلومات أصلية نقية دون أي ضغط أو معالجة. إنها مثل المكونات النيئة التي تحتاج إلى “طهي” لاحقاً لتصبح جاهزة. لأن الهواتف محدودة في التخزين وقوة المعالجة، فالصور الافتراضية هي JPEG أو HEIF “مطهية بالفعل”، ويحتاج RAW إلى تفعيل خاص.

\subsubsection*{مقارنة RAW مع HEIF و JPEG:}

\begin{itemize}
  \item \textbf{JPEG}: أقدم صيغة، أفضل توافق (أي جهاز يفتحها). عمق لون 8-bit، ضغط مع فقدان، أصغر حجم (2-5MB للصورة). مساحة محدودة للمعالجة اللاحقة، ورفع سطوح الظلام بقوة يسبب ظهور كتل لونية.
  \item \textbf{HEIF/HEIC}: صيغة جديدة تروج لها Apple، بنفس الجودة حجمها نصف JPEG (1-3MB للصورة). تدعم عمق لون 10-bit وقناة شفافية. العيب هو أن بعض الأجهزة القديمة وWindows تحتاج تثبيت إضافة لفتحها.
  \item \textbf{RAW}: بيانات أصلية غير مضغوطة، أكبر حجم (20-100MB للصورة). عمق اللون عادةً 12-16bit، أكبر مساحة للمعالجة اللاحقة. العيب هو وجوب استخدام برامج احترافية (Lightroom، Snapseed وغيرها) لفتحها ومعالجتها، ولا يمكن مشاركتها مباشرة.
\end{itemize}

\subsubsection*{لماذا يصوّر المصورون المحترفون بصيغة RAW:}

\begin{itemize}
  \item \textbf{ضبط توازن الأبيض لاحقاً بدون فقدان}: إذا أخطأت في توازن الأبيض مع JPEG (مثلاً ميل للأصفر تحت إضاءة المصباح)، فإن التصحيح اللاحق يفقد الجودة؛ أما RAW فيمكن إعادة ضبط حرارة اللون بدون أي فقدان.
  \item \textbf{تسامح كبير مع التعريض}: الظلام في JPEG المعتم بدرجتين أسود تماماً، وعند تفتيحه يظهر ضجيج كامل؛ أما RAW فيمكن "إنقاذ" تفاصيل قابلة للاستخدام من ظلام معتم بـ 3-4 درجات. وينطبق الشيء نفسه على التعريض الزائد.
  \item \textbf{تفاصيل أغنى في الظلال والإضاءات}: النطاق الديناميكي لـ RAW أكبر بكثير من JPEG، مع هامش أكبر بكثير لعدم تعرض الإضاءات الزائد وظلام الظلال التام.
  \item \textbf{الحدّة وتقليل الضجيج وتصحيح العدسة قابلة للضبط لاحقاً}: معالجة الكاميرا الداخلية "صفقة نهائية"، أما RAW فيمكن تعديله مراراً وتكراراً، انطلاقاً من البيانات الأصلية في كل مرة.
\end{itemize}

\subsubsection*{كيفية تفعيل RAW على الهاتف:}

\begin{itemize}
  \item \textbf{iPhone ProRAW}: يدعمه iPhone 12 Pro فأحدث من إصدارات Pro فقط. افتح "الإعدادات الكاميرا الصيغة" وفعّل "Apple ProRAW". سيظهر زر "RAW" في أعلى واجهة الكاميرا، اضغط لإضاءته. حجم الملف حوالي 25MB للصورة، يُنصح بتفعيله عند الحاجة فقط.
  \item \textbf{Android Motion Raw / Ultra RAW}: تختلف التسميات بين العلامات. هواتف Xiaomi/OPPO/vivo الرائدة تجد مفتاح RAW في "الوضع الاحترافي" أو "المزيد"؛ Huawei يُفعّل في "الوضع الاحترافي"؛ Samsung Galaxy S Series لديه خيار RAW في "الوضع الاحترافي". حجم الملف عادةً 20-50MB.
  \item \textbf{نصيحة عامة}: للتصوير اليومي العادي لا داعي لتفعيل RAW، فالملفات كبيرة وتشغل مساحة؛ فعّله فقط في المشاهد المهمة (المناظر الطبيعية، البورتريه، عند الحاجة للمعالجة اللاحقة).
\end{itemize}

\booksubsection{سادساً: مبادئ التصوير الحسابي}

لم يعد التصوير بالهاتف الحديث يستغني عن “التصوير الحسابي” — أي الاعتماد على الخوارزميات وليس فقط العدسات والمستشعرات لتوليد الصور. فهم هذه المبادئ الأساسية يجعلك أكثر إدراكاً لما يفعله الهاتف في لحظة الضغط على زر الغالق.

\subsubsection*{دمج الإطارات المتعددة:}

عند ضغط الغالق مرة واحدة، يلتقط الهاتف في الواقع عدة صور. تقوم الخوارزمية بمحاذاة الإطارات المتعددة ودمجها مرجّحة، للحصول على صورة عالية الجودة.

\begin{itemize}
  \item \textbf{وضع التصوير الليلي}: التقاط عدة صور معتمة في 1-3 ثوانٍ، محاذاتها وتركيبها — يلغي الضجيج المتبادل ويتراكم السطوع. والنتيجة هي الحصول على صورة نظيفة حتى في الإضاءة الخافتة.
  \item \textbf{HDR (النطاق الديناميكي العالي)}: التقاط ثلاث (أو أكثر) صور معتمة وعادية وفائقة التعريض، ودمجها في صورة واحدة تحتوي تفاصيل في الإضاءات والظلال. فعّال بشكل خاص في المشاهد المعاكسة للإضاءة.
  \item \textbf{Deep Fusion (Apple) / Photonic Engine}: مخصّص للمشاهد ذات الإضاءة المتوسطة إلى المنخفضة، يحلل الإطارات المتعددة بكسلاً ببكسل ويدمجها، مع التركيز على تحسين النسيج والتفاصيل، فتفاصيل مثل القماش المنسوج والشعر تصبح أوضح بشكل ملحوظ.
\end{itemize}

\subsubsection*{التقسيم الدلالي:}

الذكاء الاصطناعي لا يعالج الصورة كاملة بنفس الطريقة، بل “يفهم” أولاً ما في المشهد — سماء، وجه، عشب، مبنى، ملابس — ثم يطبق استراتيجيات معالجة مختلفة لكل منطقة.

\begin{itemize}
  \item \textbf{السماء}: تعزيز اللون الأزرق منفرداً وخفض السطوع لتجنب التعريض الزائد.
  \item \textbf{الوجه}: التعرف على مناطق لون البشرة، تفتيحها وتنعيمها منفرداً، دون المساس بتفاصيل الملامح والشعر.
  \item \textbf{العشب/أوراق الشجر}: تعزيز تشبع اللون الأخضر ورفع حدّة التفاصيل.
  \item \textbf{تعتيم الخلفية}: في وضع البورتريه، يتعرف الذكاء الاصطناعي على حدود الهدف، فيُصبح الخلفية ضبابية والهدف واضحاً.
\end{itemize}

\subsubsection*{تطبيقات التعلم العميق:}

\begin{itemize}
  \item \textbf{التعرف على المشاهد بالذكاء الاصطناعي}: تتعرف الكاميرا تلقائياً على مشاهد مثل "طعام" "حيوان أليف" "تصوير ليلي" "مستند" وتحسّن المعاملات وفقاً لذلك. مثلاً وضع الطعام يعزز الألوان الدافئة والتشبع ليبدو الطعام أكثر إغراءً.
  \item \textbf{تعتيم البورتريه بالذكاء الاصطناعي}: عدسات الهاتف لا تستطيع فيزيائياً تحقيق عمق ميدان ضحل كالكاميرات الاحترافية، بل تعتمد على نماذج التعلم العميق للتنبؤ بحواف الهدف ثم تعتيم الخلفية. أحدث الخوارزميات تعالج الآن تفاصيل مثل خصلات الشعر وحواف الزجاج بشكل صحيح.
  \item \textbf{إزالة بالذكاء الاصطناعي/توسيع الصورة بالذكاء الاصطناعي}: تحديد المارة أو الأشياء في المشهد، فيُكمل الذكاء الاصطناعي الأجزاء المحجوبة بناءً على السياق. هذا تطبيق للذكاء الاصطناعي التوليدي في مجال التصوير.
\end{itemize}

\subsubsection*{التصوير الحسابي مقابل التصوير البصري:}

\begin{itemize}
  \item \textbf{مزايا التصوير الحسابي}: حجم صغير مع جودة جيدة، وإمكانية التصوير في ظروف قاسية كالإضاءة الخافتة والإضاءة الخلفية، والمستخدم العادي لا يحتاج لمعرفة تصوير لالتقاط صور "جميلة".
  \item \textbf{عيوب التصوير الحسابي}: الصورة أحياناً "معالجة زيادة" — حدّة مفرطة، ألوان مشوّهة، لون سماء غير طبيعي؛ دمج الإطارات المتعددة يترك ظلاً مزدوجاً عند حركة الهدف؛ التفاصيل "محسوبة" وليست مسجلة فعلياً.
  \item \textbf{مزايا التصوير البصري}: مستشعر كبير + عدسة بفتحة واسعة تلتقط الضوء فعلياً، جودة طبيعية ومساحة كبيرة للمعالجة اللاحقة؛ لا تُولّد الخوارزميات أثراً زائفاً بسبب "التخمين".
  \item \textbf{الاتجاه}: الهواتف الرائدة الحديثة تجمع بين "البصري والحسابي" — مستشعر كبير للحصول على معلومات حقيقية، وخوارزميات لتحسين العرض. الاثنان ليسا متعارضين.
\end{itemize}

\booksubsection{سابعاً: أنواع العدسات والبعد البؤري}

الهواتف الرائدة الحديثة مزودة عادةً بثلاث إلى أربع عدسات، لكل منها بُعد بؤري مختلف. فهم البعد البؤري يجعلك تختار العدسة المناسبة في كل مشهد.

\subsubsection*{العدسات الشائعة والبعد البؤري المكافئ:}

\begin{itemize}
  \item \textbf{فائق الاتساع (13-24mm مكافئ)}: أوسع مجال رؤية، مناسب للمناظر الطبيعية والمباني والمساحات الداخلية. يلتقط مشاهد واسعة، لكن الحواف تتشوّه (الخطوط المستقيمة تنحني)، وفي تصوير البورتريه يتمدد الوجه. مناسب لـ "المشاهد الكبيرة"، غير مناسب لـ "اللقطة المقربة".
  \item \textbf{العدسة الرئيسية (24-28mm مكافئ)}: أكثر عدسة استخداماً يومياً، أفضل جودة وأكبر فتحة عدسة. توازن بين المجال والجودة، تصلح لتصوير الأشخاص والمناظر. "1x" في معظم الهواتف هي العدسة الرئيسية.
  \item \textbf{العدسة المقربة (70-200mm مكافئ)}: لتصوير المشاهد البعيدة ولقطات البورتريه المقربة. العدسة المقربة تضغط الإحساس بالعمق (الجبل البعيد والشخص القريب يبدوان أقرب)، وتعطي نسبة ملامح أكثر طبيعية وتعتيم خلفية في تصوير البورتريه.
  \item \textbf{عدسة مقربة عاكسة (200mm+ مكافئ)}: من خلال تصميم مسار ضوئي عاكس، يحقق بُعداً بؤرياً أطول ضمن سماكة الهاتف. يلتقط تفاصيل واضحة من مسافة بعيدة، مناسب للحفلات الموسيقية والفعاليات الرياضية والحياة البرية. متاح فقط في الهواتف الرائدة.
\end{itemize}

\subsubsection*{الموازنة بين حجم المستشعر وحجم البكسل:}

المستشعر هو “القاعدة”، ويحدد كمية الضوء المستقبلة. في الصناعة يُقال “قاعدة أكبر تتفوق” — بنفس عدد البكسلات، كلما كبر المستشعر، كبر كل بكسل، وزاد الضوء الداخل، وتحسنت الجودة.

\begin{itemize}
  \item \textbf{مستشعر كبير}: كمية ضوء كبيرة، جودة تصوير ليلي جيدة، نطاق ديناميكي واسع. العيب هو زيادة سماكة الهاتف (بروز وحدة العدسة) وارتفاع التكلفة.
  \item \textbf{بكسل عالي ومستشعر صغير}: نهاراً مع ضوء كافٍ، البكسل العالي يلتقط صوراً أوضح، ملائمة للقص اللاحق. لكن في الإضاءة الخافتة لا يكفي الضوء لكل بكسل، فيكثر الضجيج.
  \item \textbf{اختيار الهواتف الرائدة}: العدسة الرئيسية غالباً 1/1.3 بوصة إلى 1 بوصة (بعض الموديلات تستخدم مستشعر 1 بوصة)، بعدد بكسلات حوالي 50 مليون — وهي نقطة التوازن بين الجودة والحجم.
\end{itemize}

\subsubsection*{تقنية دمج البكسلات (Pixel Binning):}

للموازنة بين البكسل العالي نهاراً والبكسل الكبير ليلاً، يستخدم المُصنّعون “دمج البكسلات” — دمج عدة بكسلات صغيرة في بكسل كبير واحد عند الإخراج.

\begin{itemize}
  \item \textbf{دمج رباعي (Quad Bayer)}: مستشعر 50 مليون بكسل، يُخرج صوراً بدقة 12.5 مليون بكسل. كل بكسل ناتج مُركّب من 4 بكسلات متجاورة، مما يضاعف كمية الضوء الداخلة 4 مرات وجودة أنظف.
  \item \textbf{دمج تساعي (Nonacell)}: يُستخدم في مستشعرات تتجاوز 100 مليون بكسل، يُخرج حوالي 11-12 مليون بكسل. نفس المبدأ، مع كمية ضوء أكبر.
  \item \textbf{التأثير الفعلي}: في وضع التصوير الافتراضي، ما تحصل عليه هو صورة منخفضة البكسل بعد الدمج — جودة جيدة وحجم صغير. عند الحاجة لتفاصيل قصوى (مثل تصوير المستندات أو المشاهد البعيدة)، يمكن التبديل يدوياً إلى "وضع البكسل العالي" لإخراج الدقة الكاملة.
\end{itemize}

\booksubsection{ثامناً: أوضاع التصوير الخاصة}

إلى جانب التصوير اليومي، هناك عدة أوضاع خاصة في الهاتف، واستخدامها الصحيح يتيح التقاط مشاهد لا يمكن تصويرها عادةً.

\subsubsection*{الوضع البانورامي:}

\begin{itemize}
  \item أثناء التصوير، حرّك الهاتف أفقياً ببطء وثبات باتجاه السهم، لا تسرع — السرعة تسبب أخطاء في المحاذاة.
  \item ابقِ الهاتف ثابتاً عمودياً، وحرّكه أفقياً فقط. إذا اهتززت عمودياً، سيصبح الأفق متموجاً.
  \item احذر عند وجود أشخاص أو سيارات متحركة — قد يظهر "نصف جسد" أو "شبح" عند الدمج.
  \item يمكن أيضاً التصوير البانورامي عمودياً — تصوير عمودي ودمج أفقي، مناسب للمباني الشاهقة أو الشلالات.
\end{itemize}

\subsubsection*{التصوير البطيء والتصوير بتقنية الفاصل الزمني:}

\begin{itemize}
  \item \textbf{الحركة البطيئة}: التصوير بمعدل إطارات عالٍ 120fps أو 240fps، ليصبح التشغيل لقطة بطيئة. مناسب لقطرات الماء ولحظات الرياضة وقفز الحيوانات الأليفة. يجب أن تكون الإضاءة كافية، فوضع الحركة البطيئة يقلل وقت تعرّض الضوء، وجودة الإضاءة الخافتة ضعيفة.
  \item \textbf{التقاط الفاصل الزمني}: التقاط عدة صور بفواصل زمنية ودمجها في فيديو تسريع. مناسب لتصوير حركة المرور والسحب وغروب الشمس وتفتح الأزهار. مدة التصوير من دقائق إلى ساعات، ويجب تثبيت الهاتف على حامل ثلاثي، وإلا ستكون الصورة مهتزة.
  \item \textbf{نصيحة صغيرة}: عند تصوير السحب بالفاصل الزمني، ثبّت الهاتف وصوّر لمدة 15-30 دقيقة، لتحصل على فيديو مدته 10 ثوانٍ تقريباً بتأثير مذهل.
\end{itemize}

\subsubsection*{وضع التصوير النجمي/الفلكي:}

\begin{itemize}
  \item يجب الابتعاد عن تلوث الضوء في المدينة، والبحث عن ضاحية أو منطقة جبلية بلا أضواء شوارع. لا يمكن تصوير درب التبانة من المدينة.
  \item يجب استخدام حامل ثلاثي، وإلا ستكون الصورة مهتزة. يجب تثبيت الهاتف بثبات لعدة ثوانٍ إلى عشرات الثواني.
  \item يجب أن يكون الطقس صحواً بلا غيوم، والقمر ليس كبيراً جداً (البدر يحجب ضوء النجوم). الليالي المظلمة في بداية أو نهاية الشهر القمري هي الأنسب.
  \item "الوضع الليلي" في iPhone يتعرف تلقائياً على السماء المرصعة بالنجوم ويطيل وقت التعريض (حتى 30 ثانية)؛ "وضع النجوم" أو "الوضع الفلكي" في هواتف Android الرائدة يعطي نتائج مشابهة. تحلَّ بالصبر، وبعد الضغط على الغالق لا تحرّك الهاتف.
\end{itemize}

\subsubsection*{وضع التصوير المجهري:}

\begin{itemize}
  \item لتصوير الأزهار والحشرات والمجوهرات ونصوص قريبة وغيرها من الأجسام الصغيرة. يمكن التركيز على مسافة قريبة تصل إلى 1-2 سنتيمتر.
  \item بعض الهواتف لديها عدسة تصوير مجهري مخصصة (غالباً 2 أو 5 ميجابكسل بمستشعر صغير)، وجودة متوسطة؛ الهواتف الرائدة تستخدم عدسة فائقة الاتساع للتصوير المجهري أيضاً بجودة أفضل.
  \item عند التصوير عن قرب تكون العدسة قريبة جداً من الهدف وتحجب الضوء، استخدم مصباح طاولة أو كشافاً للإضاءة المساعدة عند الحاجة.
  \item مسافة التركيز قريبة جداً وعمق الميدان ضحل، وقد يفقد التركيز بتحرك الهدف أماماً أو خلفاً بضعة مليمترات. يجب أن تكون اليد ثابتة عند التصوير، والأفضل إبقاء الهدف ثابتاً.
\end{itemize}

\begin{tcbinfo}{نصائح عملية للتصوير الثابت}
\begin{itemize}
  \item التصوير اليومي: HEIF/JPEG الافتراضي كافٍ، حجم الملف صغير والجودة مقبولة.
  \item المشاهد المهمة (مناظر طبيعية، بورتريه، تحتاج معالجة لاحقة): فعّل RAW للاحتفاظ بأكبر مساحة للمعالجة اللاحقة.
  \item اختر العدسة المناسبة: المناظر الطبيعية بفائق الاتساع، البورتريه بالعدسة المقربة، والتصوير اليومي بالعدسة الرئيسية.
  \item التصوير الليلي والنجوم: ثبّت الهاتف بالضرورة، وفعّل الوضع الليلي أو الوضع الفلكي.
  \item الهواتف ذات المستشعر الكبير (مثل مستشعر 1 بوصة) لها سقف جودة أعلى، لكن عادات التشغيل تبقى عاملاً حاسماً.
\end{itemize}

\end{tcbinfo}

\begin{tcbpractice}{جرّب بنفسك: تعرّف على قدرات كاميرتك}
\begin{enumerate}
  \item تحقق مما إذا كان هاتفك يدعم التصوير بصيغة RAW (الإعدادات الكاميرا الصيغة)، حاول تفعيله والتقاط صورة، وقارن حجم ملف JPEG مع RAW لنفس المشهد.
  \item بدّل بين العدسات فائقة الاتساع والرئيسية والمقربة، وصوّر نفس المشهد البعيد بكل منها، لتدرك الفرق في زاوية الرؤية والجودة.
  \item ابحث عن مشهد يصلح للوضع الليلي أو وضع النجوم، وثبّت الهاتف ثم صوّر صورة ليلية.
  \item حاول استخدام الوضع البانورامي لتصوير مشهد واسع أفقياً، مع الحفاظ على تحريك الهاتف بسلاسة.
\end{enumerate}

\end{tcbpractice}
